% GENERATED FILE. Edit canonical lessons, then run npm run generate:books.
% canonical-source-hash: fnv1a64:3846115ad02fac49
% canonical-lessons: MR-C412-disa, MR-C412-hokayantra, MR-C412-khun, MR-C412-bharkatne, MR-C412-ghutmalne

\chapter{Finding the Way}
\label{ch:mr-finding-the-way}
\hlchaptermodality{412}

\begin{chapteropening}
\textbf{By the end of this chapter:} \emph{I can talk about directions, a compass, landmarks, and getting lost.}

Complete the last lesson of chapter 412: I can talk about directions, a compass, landmarks, and getting lost.
\end{chapteropening}

\section[\texorpdfstring{\mr{दिशा} (diśā)}{दिशा (diśā)}]{\mr{दिशा} (diśā) — a direction}

\label{lesson:MR-C412-disa}



\begin{quote}
Before the new one: say the Marathi for to roam, then the Marathi for to linger.
\end{quote}

\subsection*{You'll want to know: \mr{दिशा}}
\textbf{\mr{दिशा}} — \emph{diśā} — \textquotedblleft{}a direction\textquotedblright{}.

\begin{quote}
\mr{पुणे} \mr{कोणत्या} \mr{दिशेला} \mr{आहे}?
\end{quote}

\begin{quote}
\emph{puṇe koṇatyā diśelā āhe?}
\end{quote}

\begin{quote}
In which direction is Pune?
\end{quote}

Feminine.

Before an ending \mr{दिशा} becomes \mr{दिशे}-: \mr{दिशेला}, \mr{दिशेने}.

\begin{cousinweb}
\textbf{\mr{दिशा}} (\emph{diśā}) is a tatsama, taken from Sanskrit \emph{diśā}, \textquotedblleft{}direction, quarter of the compass\textquotedblright{} (Monier-Williams).
\end{cousinweb}

\begin{grammarlens}[title={What you've built}]
One more word for getting around.
\end{grammarlens}

\subsection*{Your turn}
\begin{itemize}
\raggedright
  \item \emph{Say it:} \emph{diśā}
  \item \emph{Say it:} \emph{diśā}, once more
  \item \emph{Say it:} \emph{puṇe koṇatyā diśelā āhe?} — In which direction is Pune?
  \item \emph{Recall:} say the Marathi for to roam, then the Marathi for to linger, then say \emph{diśā} again
\end{itemize}

\begin{tcolorbox}[breakable,colback=teal!4,colframe=teal!35!black,title={Before you move on}]
What does \mr{दिशा} mean? (\textquotedblleft{}A direction\textquotedblright{}.) Say it once more.
\end{tcolorbox}
\section[\texorpdfstring{\mr{होकायंत्र} (hokāyantra)}{होकायंत्र (hokāyantra)}]{\mr{होकायंत्र} (hokāyantra) — a compass}

\label{lesson:MR-C412-hokayantra}



\begin{quote}
Before the new one: say the Marathi for to linger, then the Marathi for a direction.
\end{quote}

\subsection*{You'll want to know: \mr{होकायंत्र}}
\textbf{\mr{होकायंत्र}} — \emph{hokāyantra} — \textquotedblleft{}a compass\textquotedblright{}.

\begin{quote}
\mr{जंगलात} \mr{होकायंत्र} \mr{उपयोगी} \mr{पडतं}.
\end{quote}

\begin{quote}
\emph{jaṅgalāt hokāyantra upyogī paḍtaṁ.}
\end{quote}

\begin{quote}
In the forest a compass comes in useful.
\end{quote}

Neuter.

The needle of a compass points north.

\begin{cousinweb}
\textbf{\mr{होकायंत्र}} (\emph{hokāyantra}) ends in \textbf{\mr{यंत्र}}, \textquotedblleft{}a device\textquotedblright{}, a tatsama from Sanskrit \emph{yantra} (Monier-Williams); where the first part comes from is not settled.
\end{cousinweb}

\begin{grammarlens}[title={What you've built}]
One more word for getting around.
\end{grammarlens}

\subsection*{Your turn}
\begin{itemize}
\raggedright
  \item \emph{Say it:} \emph{hokāyantra}
  \item \emph{Say it:} \emph{hokāyantra}, once more
  \item \emph{Say it:} \emph{jaṅgalāt hokāyantra upyogī paḍtaṁ.} — In the forest a compass comes in useful.
  \item \emph{Recall:} say the Marathi for to linger, then the Marathi for a direction, then say \emph{hokāyantra} again
\end{itemize}

\begin{tcolorbox}[breakable,colback=teal!4,colframe=teal!35!black,title={Before you move on}]
What does \mr{होकायंत्र} mean? (\textquotedblleft{}A compass\textquotedblright{}.) Say it once more.
\end{tcolorbox}
\section[\texorpdfstring{\mr{खूण} (khūṇ)}{खूण (khūṇ)}]{\mr{खूण} (khūṇ) — a mark; a landmark}

\label{lesson:MR-C412-khun}



\begin{quote}
Before the new one: say the Marathi for a direction, then the Marathi for a compass.
\end{quote}

\subsection*{You'll want to know: \mr{खूण}}
\textbf{\mr{खूण}} — \emph{khūṇ} — \textquotedblleft{}a mark; a landmark\textquotedblright{}.

\begin{quote}
\mr{मंदिर} \mr{ही} \mr{आमच्या} \mr{गल्लीची} \mr{खूण} \mr{आहे}.
\end{quote}

\begin{quote}
\emph{mandir hī āmcyā gallīcī khūṇ āhe.}
\end{quote}

\begin{quote}
The temple is the landmark of our lane.
\end{quote}

Feminine.

\mr{खूण} \mr{करणे}, to make a mark or a sign with the hand.

\begin{cousinweb}
\textbf{\mr{खूण}} (\emph{khūṇ}) is a Marathi word for a mark, a sign or a landmark; its origin is not settled.
\end{cousinweb}

\begin{grammarlens}[title={What you've built}]
One more word for getting around.
\end{grammarlens}

\subsection*{Your turn}
\begin{itemize}
\raggedright
  \item \emph{Say it:} \emph{khūṇ}
  \item \emph{Say it:} \emph{khūṇ}, once more
  \item \emph{Say it:} \emph{mandir hī āmcyā gallīcī khūṇ āhe.} — The temple is the landmark of our lane.
  \item \emph{Recall:} say the Marathi for a direction, then the Marathi for a compass, then say \emph{khūṇ} again
\end{itemize}

\begin{tcolorbox}[breakable,colback=teal!4,colframe=teal!35!black,title={Before you move on}]
What does \mr{खूण} mean? (\textquotedblleft{}A mark; a landmark\textquotedblright{}.) Say it once more.
\end{tcolorbox}
\section[\texorpdfstring{\mr{भरकटणे} (bharkaṭṇe)}{भरकटणे (bharkaṭṇe)}]{\mr{भरकटणे} (bharkaṭṇe) — to stray; to drift}

\label{lesson:MR-C412-bharkatne}



\begin{quote}
Before the new one: say the Marathi for a compass, then the Marathi for a mark.
\end{quote}

\subsection*{You'll want to know: \mr{भरकटणे}}
\textbf{\mr{भरकटणे}} — \emph{bharkaṭṇe} — \textquotedblleft{}to stray; to drift\textquotedblright{}.

\begin{quote}
\mr{धुक्यात} \mr{आम्ही} \mr{वाटेवरून} \mr{भरकटलो}.
\end{quote}

\begin{quote}
\emph{dhukyāt āmhī vāṭevarūn bharkaṭlo.}
\end{quote}

\begin{quote}
In the fog we strayed from the path.
\end{quote}

Past (he …): \mr{तो} \mr{भरकटला}, \emph{to bharkaṭlā}.

For people who lose their way, and for a mind that wanders.

\begin{cousinweb}
\textbf{\mr{भरकटणे}} (\emph{bharkaṭṇe}) is a Marathi verb for drifting off the right path; its origin is not settled.
\end{cousinweb}

\begin{grammarlens}[title={What you've built}]
One more word for getting around.
\end{grammarlens}

\subsection*{Your turn}
\begin{itemize}
\raggedright
  \item \emph{Say it:} \emph{bharkaṭṇe}
  \item \emph{Say it:} \emph{bharkaṭṇe}, once more
  \item \emph{Say it:} \emph{dhukyāt āmhī vāṭevarūn bharkaṭlo.} — In the fog we strayed from the path.
  \item \emph{Recall:} say the Marathi for a compass, then the Marathi for a mark, then say \emph{bharkaṭṇe} again
\end{itemize}

\begin{tcolorbox}[breakable,colback=teal!4,colframe=teal!35!black,title={Before you move on}]
What does \mr{भरकटणे} mean? (\textquotedblleft{}To stray; to drift\textquotedblright{}.) Say it once more.
\end{tcolorbox}
\section[\texorpdfstring{\mr{घुटमळणे} (ghuṭmaḷṇe)}{घुटमळणे (ghuṭmaḷṇe)}]{\mr{घुटमळणे} (ghuṭmaḷṇe) — to hang about; to hesitate}

\label{lesson:MR-C412-ghutmalne}



\begin{quote}
Before the new one: say the Marathi for a mark, then the Marathi for to stray.
\end{quote}

\subsection*{You'll want to know: \mr{घुटमळणे}}
\textbf{\mr{घुटमळणे}} — \emph{ghuṭmaḷṇe} — \textquotedblleft{}to hang about; to hesitate\textquotedblright{}.

\begin{quote}
\mr{तो} \mr{दुकानाबाहेर} \mr{बराच} \mr{वेळ} \mr{घुटमळला}.
\end{quote}

\begin{quote}
\emph{to dukānābāher barāc veḷ ghuṭmaḷlā.}
\end{quote}

\begin{quote}
He hung about outside the shop for quite a while.
\end{quote}

Past (he …): \mr{तो} \mr{घुटमळला}, \emph{to ghuṭmaḷlā}.

Moving a little but not leaving, as if unsure.

\begin{cousinweb}
\textbf{\mr{घुटमळणे}} (\emph{ghuṭmaḷṇe}) is a Marathi verb for hovering in one place, unsure whether to go; its origin is not settled.
\end{cousinweb}

\begin{grammarlens}[title={What you've built}]
One more word for getting around.
\end{grammarlens}

\subsection*{Your turn}
\begin{itemize}
\raggedright
  \item \emph{Say it:} \emph{ghuṭmaḷṇe}
  \item \emph{Say it:} \emph{ghuṭmaḷṇe}, once more
  \item \emph{Say it:} \emph{to dukānābāher barāc veḷ ghuṭmaḷlā.} — He hung about outside the shop for quite a while.
  \item \emph{Recall:} say the Marathi for a mark, then the Marathi for to stray, then say \emph{ghuṭmaḷṇe} again
\end{itemize}

\begin{tcolorbox}[breakable,colback=teal!4,colframe=teal!35!black,title={Before you move on}]
What does \mr{घुटमळणे} mean? (\textquotedblleft{}To hang about; to hesitate\textquotedblright{}.) Say it once more.
\end{tcolorbox}
